% Generated draft: requires mathematical and source review.
\documentclass[11pt]{article}
\usepackage[margin=25mm]{geometry}
\usepackage{fontspec}
\setmainfont{Latin Modern Roman}
\usepackage{amsmath,amssymb,mathtools,unicode-math}
\setmathfont{Latin Modern Math}
\usepackage{enumitem,needspace,xurl,hyperref}
\setlength{\parindent}{0pt}
\setlength{\parskip}{5pt}
\setlength{\emergencystretch}{3em}
\newcommand{\N}{\mathbb{N}}
\newcommand{\Z}{\mathbb{Z}}
\newcommand{\Q}{\mathbb{Q}}
\newcommand{\R}{\mathbb{R}}
\newcommand{\C}{\mathbb{C}}
\newcommand{\F}{\mathbb{F}}
\newcommand{\A}{\mathbb{A}}
\newcommand{\PP}{\mathbb{P}}
\newcommand{\E}{\mathbb{E}}
\newcommand{\eps}{\varepsilon}
\newcommand{\dd}{\,\mathrm{d}}
\begin{document}
\section*{UnsolvedMath ID: 800465}
\textbf{Source problem code:} catalogue entry
\subsection{Definitions and mathematical statement}
Let \((M^n,g)\) be an \(n\)-dimensional Riemannian manifold and let \(p\in M\). Write
\[
B_1^g(p)=\{x\in M:d_g(x,p)<1\}
\]
for the open geodesic unit ball, and let \(\dd\operatorname{vol}_g\) denote Riemannian volume measure. In the standard interpretation of the supplied notation, \(R=\operatorname{Scal}_g=\operatorname{tr}_g\operatorname{Ric}_g\) is scalar curvature, and the lower Ricci bound means
\[
\operatorname{Ric}_g(v,v)\geq-\lambda g(v,v)
\qquad\text{for every }x\in M\text{ and }v\in T_xM.
\]
The source presents the following local version of a conjecture attributed to Yau: there exists a finite constant \(C(n,\lambda)\), depending only on the dimension and the Ricci lower-bound parameter, such that every manifold and every center satisfying the intended hypotheses obey
\[
\operatorname{Ric}_g\geq-\lambda g
\quad\Longrightarrow\quad
\int_{B_1^g(p)}\operatorname{Scal}_g\,\dd\operatorname{vol}_g
\leq C(n,\lambda).
\]
The integral is unnormalized and has no absolute value. The surrounding source discussion explicitly includes the possibility of small volume; no positive lower bound on the ball volume is stated.

The supplied statement does not specify completeness, boundary conditions, the admissible range of \(\lambda\), or any additional restrictions on the unit ball. These qualifications remain unresolved and have not been inserted as hypotheses. The interpretations of \(R\), the tensor inequality, and the integration measure above are standard readings of notation that the supplied excerpt does not explicitly define. Consequently, this is a statement-only rendering of the supplied conjecture with an incompletely specified class of admissible manifolds.
\subsection{Short English statement}
The local conjecture attributed to Yau asks whether a lower Ricci curvature bound controls the total scalar curvature of every geodesic unit ball by a constant depending only on \(n\) and \(\lambda\), including in situations with small volume. The supplied excerpt leaves completeness, boundary conditions, and the range of \(\lambda\) unspecified.
\subsection{Sources}
\begin{itemize}
\item Supplied local UnsolvedMath registry record: permanent ID \(800465\), exact source problem code \texttt{catalogue entry}, title ``Spaces with Lower Ricci Curvature --- Conjecture 9 (local Yau conjecture).'' The registry labels the record open; that label is supplied metadata, not an independently verified mathematical or literature assessment.
\item Supplied source attribution: Naber, \emph{Conjectures and Open Questions on Spaces with Lower Ricci Curvature} (2020), Conjecture 9, \url{https://arxiv.org/abs/2010.10031}. The record describes extraction from arXiv LaTeX and reports access on 2026-10-16; these provenance details are supplied and have not been independently verified here.
\item Attribution qualification: the registry lists Naber as proposer, while the supplied statement attributes the main conjecture to Yau and presents this formulation as a local version.
\end{itemize}
\end{document}
